\section{Validation: from RMSE to decision-aware metrics}
\label{sec:validation}

Most surrogate-modeling studies in energy-system optimization use RMSE,
MAE and $R^{2}$ as validation metrics for a predictive model. However,
whether the solution found with the surrogate is feasible in the actual
optimization model is then treated merely as an application of the
already \"validated\" model, without specifically testing the surrogate
in the optimizer.

This section argues that the dominant validation practice is necessary
but not always sufficient. A surrogate that scores well on a regression
benchmark can still produce solutions that are infeasible, dominated or
unstable when used inside an optimizer, because the regression metric
averages errors uniformly while the optimizer reacts to errors in active
constraints and in regions where the posterior -- the predicted mean and
spread at each input -- is locally biased.

The subsections below review validation practice: standard point and
interval metrics (SectionÂ~8.1); feasibility and violation severity
(SectionÂ~8.2); decision-aware cost and gap metrics
(SectionÂ~8.3); stress tests under extreme operating
conditions (SectionÂ~8.4); out-of-distribution and policy-shift
evaluation (SectionÂ~8.5); and reproducibility of reported results
(SectionÂ~8.6). The corresponding metrics and test designs
are summarized in
TableÂ~\ref{tab:T4-validation}.

\begin{table*}[!t]
\centering
\scriptsize
\begin{tabularx}{\textwidth}{@{}L{0.22\textwidth} L{0.16\textwidth} L{0.26\textwidth} L{0.28@\textwidth} Metric / test & Tells you about & When it is@{}}
\toprule
useful (refs.) & Common misuses (refs.)\
\\
RMSE / MAE / $R^2$ & Average fit & Sanity / first screen
\cite{Safta20172535,Schulte2020,Chen2022,Xu20223066} &
Reported alone, without a decision-aware companion\
Quantile / prediction-interval coverage & Calibration of posterior
spread & Probabilistic surrogates, BO acquisition
\cite{Xu20212696,Hu2024,Buchanan2024} &
Optimistic on the training distribution; calibration not checked after
distribution shift\
\\
Decision regret (cost gap to oracle) & Loss in decision quality &
Surrogate inside optimizer (P1, P3)
\cite{Chen20244723,Wu2022231,Liang20246508,Wu2024391} & Mean gap
reported without its distribution or a fixed solver-time comparison\
Feasibility rate & Constraint safety on the host problem &
Hard-constrained problems (P1, P3, P5)
\cite{Chen2022,Chen20244723,Liu20223726,Liang20246508} & Reported
without violation severity\
Violation severity (worst case + 95th percentile) & Robustness under
tail load / contingencies & Safety-critical operation, stress regimes
\cite{Chen2023657,Liang20246508,Chen20244723} & Average-only reporting;
missing contingency stratification\
Objective gap distribution (MIP gap, CVaR of gap) & Distributional
fairness across instances & MILP / MINLP surrogate use
\cite{Wu2022231,Chen20244723,Wu2024391} & Mean-only summaries hide tail
behaviour\
Distribution-shift test & Robustness beyond the training distribution &
Changing topology, capacity mix or operating conditions
\cite{Park2019,Subedi2026157,Liu20222679} & A single random IID split\
\\
Stress tests (peak load, low-renewable spell, cold spell) & Tail
behaviour & ESM operation under extremes
\cite{Hu2024,Buchanan2024,Bayani20231024,Sun20191497} & Only on
average days; no explicit extreme-event block\
Counterfactual policy shift (capacity mix, tariff regime) & Robustness
\cite{Park2019} & Validation only on the historical or baseline mix\
Scenario-shift / distribution-shift split & Generalisation beyond IID &
Stochastic / robust use, OOD planning
\cite{Subedi2026157,Liu20222679} & IID random split only\
Holdout reproducibility (seeds, splits, code, data) & Trust and
comparability & All surrogate classes \cite{Khazeiynasab2021,Zhang2025__7}
& Single seed or split; missing code or data release\
\end{tabularx}
\end{table*}

\subsection{Point and interval prediction metrics}
\label{sec:val:point}

Point and interval metrics remain the most widely employed validation
metrics. Given a held-out test set of $n$ pairs $(\mathbf{x}_i, y_i)$,
where $y_i$ denotes a quantity returned by the high-fidelity model and
$\hat{y}_i$ the surrogate prediction, the most common point metrics are
$$\begin{align}
  \mathrm{MAE}  &= \frac{1}{n}\sum_{i=1}^{n}\lvert y_i - \hat{y}_i\rvert,
  \label{eq:mae} \\
  \mathrm{RMSE} &= \sqrt{\frac{1}{n}\sum_{i=1}^{n}(y_i - \hat{y}_i)^2},
  \label{eq:rmse} \\
  R^{2}         &= 1 - \frac{\sum_{i=1}^{n}(y_i - \hat{y}_i)^2}
                       {\sum_{i=1}^{n}(y_i - \bar{y})^2},
  \label{eq:r2}
\end{align}$$

with $\bar{y} = \frac{1}{n}\sum_i y_i$. MAE averages absolute errors and
is easy to interpret in the units of $y$; RMSE penalizes large errors
more heavily; $R^{2}$ measures how much variance in $y$ is explained
relative to a constant baseline. Studies that employed these metrics in
the energy system optimization domain include, for example, polynomial
response surfaces for integrated electric--thermal scheduling
\cite{Wang2020202933}), regional integrated energy system loads
\cite{Shi2024__2}), LSTM surrogates for DH with thermal storage \cite{Du2025}),
GP surrogates in MES microgrids \cite{Li2023__3}), and NN surrogates for
distributed energy-system and MES optimization
\cite{Perera2019191,Ren2024}).

For probabilistic surrogates, validation additionally requires that the
reported uncertainty band matches empirical coverage. A common check is
the prediction-interval coverage probability (PICP): the fraction of
test points that fall inside a $(1-\alpha)$ prediction interval
$[\hat{y}_i^{\mathrm{lo}},
\hat{y}_i^{\mathrm{hi}}]$,

$$\begin{equation}
  \mathrm{PICP} = \frac{1}{n}\sum_{i=1}^{n}
  \mathbb{I}\!\left[\hat{y}_i^{\mathrm{lo}} \le y_i \le
  \hat{y}_i^{\mathrm{hi}}\right],
  \label{eq:picp}
\end{equation}$$

where $\mathbb{I}[\cdot]$ is the indicator function. Pinball loss and
the continuous ranked probability score (CRPS) summarize quantile and
full predictive-distribution accuracy, respectively
\cite{Hu2024,Sun20191497}. Representative interval checks include
Transformer--GP net-load forecasts \cite{Hu2024}), GPR renewable intervals
for microgrid robust scheduling \cite{Yang2024__2}), linked GP emulators in
municipal energy planning \cite{Volodina2022,Lawson201647}, and GP
emulation of stochastic dispatch under wind uncertainty
\cite{Hu2020,Hu2021671}. CNN-based inverter surrogates can retain low bulk
RMSE while interval calibration breaks down in operating regions absent
from the training set \cite{Subedi2026157}. The literature therefore treats
point and interval scores as a screening filter rather than as a final
acceptance criterion for surrogate- assisted optimization
\cite{Khaloie2025,Lim2025}.

\subsection{Constraint feasibility and violation severity}
\label{sec:val:feasibility}

When the surrogate is used inside the optimizer (P1, P3, P5),
feasibility matters more than the point and interval prediction metrics
described above because the key question is not how well it predicts on
average, but whether the operating plan or design it produces is
actually feasible. The primary validation metric is, hence, the
feasibility rate $\phi_{\mathrm{feas}}$
(Eq.Â~\eqref{eq:feas-rate}): on many test situations that were excluded
from training, a plan is generated for each situation by the
surrogate-based optimizer, and the share of plans that satisfy all
limits of the full model (for example generation caps, line flows or
voltage bounds) is recorded.

$$\begin{equation}
  \phi_{\mathrm{feas}}
    = \frac{1}{N}\sum_{j=1}^{N}
      \mathbb{I}\!\big[\mathbf{x}_j \in \mathcal{F}\big],
  \label{eq:feas-rate}
\end{equation}$$

where $\mathcal{F}$ is the set of decisions that satisfy every
constraint of the full optimization model. When a decision is
infeasible, it also matters how far it exceeds the limits. The largest
breach across the $N$ scenarios and the size of typical large breaches
(for example at the 95th percentile) indicate whether errors remain
small or become critical under e.g. peak load. When a constraint must
hold only with a stated probability (chance constraints), the same check
is applied on hold-out scenarios: whether violations occur at most at
the promised rate, not only whether the surrogate error is small on
average.

For unit commitment, learned warm starts are evaluated for constraint
satisfaction on line losses and unit limits \cite{Chen2022}. In hybrid AC/DC
microgrid energy management, hard constraints are enforced during
training of ReLU neural surrogates for converter losses, and feasibility
of the resulting stochastic unit- commitment solutions is assessed
\cite{Liu20242534}. For joint chance-constrained OPF with renewable
generation, feasibility of quantile-regression surrogates is checked
against chance limits \cite{Chen2023657}. For PV hosting-capacity planning
in distribution networks, violations of voltage and thermal limits are
reported relative to hosting limits \cite{Koirala20242284}.

\subsection{Decision-aware metrics: regret, cost gap, MIP-gap distribution}
\label{sec:val:decision}

When the surrogate sits inside the optimizer, the validation question
shifts from prediction accuracy
(SectionÂ~8.1) and constraint satisfaction
(SectionÂ~8.2) to objective quality: how much extra
cost or lost revenue is caused by the surrogate-based model when
evaluated in the high-fidelity model. As in
SectionÂ~8.2, each test scenario is one operating
condition excluded from surrogate training. On that scenario, the
surrogate-based optimizer returns an operating decision
$\hat{\mathbf{x}}$ and the high-fidelity model yields a reference
optimum $\mathbf{x}^{*}$. The most common scalar summary is the relative
cost gap (also called regret or optimality gap),

Eq.Â~\eqref{eq:regret}: $$\begin{equation}
  r \;=\; \frac{f(\hat{\mathbf{x}}) - f(\mathbf{x}^{*})}{f(\mathbf{x}^{*})},
  \label{eq:regret}
\end{equation}$$

where $f(\cdot)$ is the objective of the high-fidelity optimization.
Solver-time reduction is usually reported alongside $r$ because a
speed-up without a bounded gap is not sufficient for acceptance. For
mixed-integer problems, studies also report the MIP gap distribution
across test scenarios under a common time limit: the share of instances
solved to optimality and how large the remaining gaps are in the slow
tail.

For large-scale economic dispatch, optimality gaps and run-time
reductions are reported for end-to-end neural proxies with repair layers
\cite{Chen20244723}. For security-constrained unit commitment with wind and
battery storage, optimality gaps are reported under a fixed solver
budget \cite{Wu2022231}. For joint chance-constrained unit commitment, a
regret percentage quantifies the extra cost of the accelerated method
relative to a reference solver \cite{Liang20246508}. In MES with a ML market
surrogate, annual revenue error relative to a validation simulation is
reported for price-taker and market-aware formulations \cite{Jalving2023}.
For unit commitment, near-optimality of overall solutions and average
gaps relative to branch-and-cut benchmarks are reported when ML
accelerates subproblem solving inside surrogate Lagrangian relaxation
\cite{Wu2024391}. Recent reviews note that decision-aware metrics are still
reported inconsistently and recommend cost-gap-style summaries as a
default \cite{Ruan2021221,Khaloie2025,Lim2025}.

\subsection{Stress tests: extreme weather, peak load, low-renewable spells}
\label{sec:val:stress}

Decision-aware metrics from
SectionsÂ~8.2
andÂ~8.3, when averaged over a uniform test set,
underweight the operating regimes that drive cost spikes and constraint
violations. A stress-test is therefore a deliberate subset of held-out
test scenarios---for example high-demand hours, cold or heat waves or
extreme weather conditions---on which the same feasibility, cost-gap or
forecast metrics are reported separately rather than integrated into the
overall mean.

Hu et al.Â~combine a Transformer network with GP residual learning and
report better net-load forecasts at peaks and during large PV
variability on a hold-out period \cite{Hu2024}. Sun et al.Â~optimize Laplace
predictive distributions with pinball loss across 99 wind-power
quantiles and evaluate nine prediction-interval coverage levels on NREL
test sites \cite{Sun20191497}. Du et al.Â~train an LSTM surrogate for an DH
system with phase-change thermal storage and report peak heating
demand---defined as the average of the three highest hourly
loads---together with price-responsive heating costs under a
peak-shaving strategy \cite{Du2025}. Yang et al.Â~use GP renewable intervals
to build robust uncertainty sets and test day-ahead to intra-day
microgrid scheduling under worst-case renewable scenarios
\cite{Yang2024__2}. Hu et al.Â~emulate stochastic economic dispatch with a
Karhunen--Loève-reduced GP and Chen et al.Â~train deep
quantile-regression surrogates for joint chance-constrained OPF and
evaluate constraint-violation quantiles under renewable uncertainty
\cite{Chen2023657}. Liang et al.Â~report violation rates when sample-based
chance constraints in accelerated unit commitment are stressed by many
extreme renewable scenarios \cite{Liang20246508}. Khaloie et al., Lim et
al.Â~note the absence of standardised stress-test
designs for surrogate-assisted energy-system optimization
\cite{Khaloie2025,Lim2025}.

\subsection{Out-of-distribution and policy-shift checks}
\label{sec:val:ood}

A random hold-out split from the same historical regime tests
interpolation within the training distribution, not extrapolation to the
optimization problem. An out-of-distribution check therefore evaluates
the surrogate on inputs deliberately drawn from a different regime than
the training data, for example a future year, climate zone, DH network
or power grid topology. A policy-shift block is the planning
counterpart: the test scenarios use a counterfactual policy or
investment decision (new storage, heat pumps, electrolyzers, tariff
rules, or market formulations) rather than only resampling past weather.

Park et al.Â~propose a physics-induced GNN for wind-farm power
estimation, report transferable performance across turbine layouts and
wind conditions, and test it in a layout-optimization experiment
\cite{Park2019}. Subedi et al.Â~integrate CNN-based inverter surrogates into
the open-source solver GridLAB-D, validate them dynamically against
reference models, and assess cross-hardware generalization with transfer
learning between inverter platforms \cite{Subedi2026157}. Volodina et
al.Â~link GP emulators across a network of municipal energy models and
propagate code uncertainty to input combinations in an OSeMOSYS-based
case study \cite{Volodina2022}. Lawson et al.Â~use statistical emulators in a
Bayesian power-network planning framework to quantify model-output
uncertainty at parameter settings that were not included in the original
design-of-experiments \cite{Lawson201647}. Liu et al.Â~design GP kernels for
data-driven probabilistic load flow and select hyperparameters with a
potential-game formulation so that the surrogate generalizes to datasets
not used in training \cite{Liu20222679}. Jalving et al.Â~compare price-taker
and market-aware energy system designs optimized with ML market
surrogates against a full validation simulation in IDAES and report much
lower annual revenue error for the market-aware formulation
\cite{Jalving2023}. Khaloie et al., Lim et al.Â~and Mylonopoulos et
al.Â~recommend that OOD and policy-shift evidence should
accompany surrogate-assisted planning studies
\cite{Khaloie2025,Lim2025,Mylonopoulos202332697}.

\subsection{Reproducibility: seeds, splits, code, data}
\label{sec:val:repro}

Reported speed-ups, regret values and feasibility rates are only
comparable when the full pipeline can be rerun. Reproducibility in
surrogate-assisted optimization therefore requires, at minimum 1) the
random seeds used to generate design points and training batches, 2) an
explicit train/validation/test split rule, 3) the surrogate training
code, 4) the optimization model and solver settings, 5) and access to
the underlying data or and code. Missing any one element makes a
decision-aware metric hard to interpret or compare across studies.

Representative positive practice includes generator-parameter
calibration with a stochastic RBF surrogate and documented calibration
objectives \cite{Khazeiynasab2021}, and fuel-cell system calibration with an
adaptive response surface, LHS and stated calibration targets
\cite{Zhang2025__7}. Recent reviews of learning-assisted power-system
optimization note uneven release of code and data, incomplete reporting
of splits and seeds, and inconsistent use of decision-aware metrics,
which limits cross- study comparison
\cite{Ruan2021221,Khaloie2025,Lim2025,Mylonopoulos202332697}. One
review explicitly lists reproducibility, open-source code and benchmark
comparisons among desirable reporting practices for learning-based OPF
\cite{Khaloie2025}.
